\documentclass[12pt]{article}
\usepackage[T1]{fontenc}
\usepackage[utf8]{inputenc}
\usepackage{lmodern}
\usepackage{textcomp}
\usepackage{enumitem}
\usepackage{geometry}
\geometry{margin=1in}
\begin{document}
\begin{center}
Answer Key - History - Graduate Level Assessment 1 \
\small (Answers are ordered exactly as in the question paper.)
\end{center}

\section*{Multiple Choice}
\begin{enumerate}[label=\arabic*.]
\item \textbf{Answer:} A
\\ \textit{Options:} A) Aztecs.; B) Greeks.; C) Sumerians.; D) Chavin.; E) Chinese.
\\ \textit{Note:} The statement is incorrect; the Maya developed their own distinct writing and mathematical systems independently, and while the Aztecs also had their own systems, the Maya were not derived from them.
\item \textbf{Answer:} B
\\ \textit{Options:} A) the Incas were suffering from a plague when Pizarro arrived.; B) the empire had never won the allegiance of the peoples it had conquered.; C) the Incas were already weakened by a civil war.; D) the Incas were expecting Pizarro and willingly surrendered to him.; E) he arrived with a superior force of more than 80,000 soldiers.
\\ \textit{Note:} The correct answer highlights that the Inca Empire's reliance on coercive control over subjugated peoples, rather than establishing genuine loyalty or allegiance, left these groups vulnerable to Pizarro's conquest, as they were less inclined to resist an external force that promised them greater autonomy or relief from Inca oppression.
\item \textbf{Answer:} A
\\ \textit{Options:} A) the construction of canals for irrigation; B) the creation of a trade network for llama products; C) the need to collect large quantities of shellfish; D) the building of the Piramide Mayor; E) the need for a centralized defense against invaders
\\ \textit{Note:} The construction of canals for irrigation at Caral facilitated agricultural surplus, which enabled the development of a complex society through increased population density, social stratification, and the establishment of trade networks.
\item \textbf{Answer:} E
\\ \textit{Options:} A) a shared ability to use complex tools; B) a similar diet; C) a reliance on brachiation; D) an identical social structure; E) a mode of locomotion
\\ \textit{Note:} Homo sapiens and Australopithecus afarensis share a mode of locomotion characterized by bipedalism, which is a defining trait of both species and highlights their evolutionary adaptation to walking on two legs.
\item \textbf{Answer:} A
\\ \textit{Options:} A) soon after 60,000 years ago; B) approximately 50,000 years ago; C) no earlier than 100,000 years ago; D) soon after 30,000 years ago; E) no later than 70,000 years ago
\\ \textit{Note:} Most archaeologists agree that the occupation of Australia began soon after 60,000 years ago due to substantial archaeological evidence, such as the discovery of ancient tools and human remains, indicating a sustained human presence in the region during that timeframe.
\end{enumerate}

\section*{True / False}
\begin{enumerate}[label=\arabic*.]
\setcounter{enumi}{5}
\item \textbf{Answer:} True
\\ \textit{Options:} A) True; B) False
\\ \textit{Note:} According to the passage: Byzantine Empire
\item \textbf{Answer:} False
\\ \textit{Options:} A) True; B) False
\\ \textit{Note:} The correct answer is: Igor Mandić
\item \textbf{Answer:} True
\\ \textit{Options:} A) True; B) False
\\ \textit{Note:} According to the passage: 1980 s
\item \textbf{Answer:} True
\\ \textit{Options:} A) True; B) False
\\ \textit{Note:} According to the passage: Manuel Doblado
\item \textbf{Answer:} True
\\ \textit{Options:} A) True; B) False
\\ \textit{Note:} According to the passage: Royal Australian Air Force
\end{enumerate}

\section*{Long Form Response}
\begin{enumerate}[label=\arabic*.]
\setcounter{enumi}{10}
\item \textbf{Answer:} Baden
\\ \textit{Note:} Expected answer: Baden
\item \textbf{Answer:} Archbishop of Canterbury Thomas Cranmer
\\ \textit{Note:} Expected answer: Archbishop of Canterbury Thomas Cranmer
\end{enumerate}

\vfill
\noindent\textit{End of Answer Key}
\end{document}